\section{From Average Effects To Individual
  Effects}\label{sec:introduction}

Estimating the average treatment effect (ATE) for a population of
interest has been a main focus of a rich literature on causal
inference, and the last decades have seen the development of extensive
statistical theory addressing important issues such as identification,
estimation, and uncertainty quantification
\citep[e.g.][]{rubin1974estimating, pearl1995causal}. That said, the
average effect only provides a coarse summary of the distribution of a
treatment effect, and may be insufficient, or even misleading, to
validate an intervention. Imagine, for instance, that a drug cures
70\% of the patients while it makes the symptoms much worse for the remaing
30\%. It is unclear whether the drug should be approved in spite of the
positive average effect. This is not an artificial example; 
trials typically do not give clear cut results and examples of this kind are common.
% trials
% typically do not come out black and white and examples of this kind
% are
% frequent. 
% \cite{kent2010assessing} present quite a few examples in medical science For instance, Tissue plasminogen activator, in the treatment of ST-Elevation acute myocardial infarction, significantly improves survival in high-risk patients compared to streptokinase while increase the excessive risk in low-risk patients which may outweigh the benefits.
On May 31, 2018, the National Academy of Medicine (NAM) held a
workshop to discuss approaches of examining individual treatment
effects (ITE) to support individualized patient care. In conjunction
with this event,  NAM published a report highlighting the importance
of ITE in medicine.  We quote from this document:
\begin{quote}
``The individuality of the patient should be at the core of every treatment decision. One-size-fits-all approaches to treating medical conditions are inadequate; instead, treatments should be tailored to individuals based on heterogeneity of clinical characteristics and their personal preferences.''
\end{quote}
Outside of medical science, treatment effect heterogeneity is also of
great concern to political scientists \citep{imai2011estimation,
  grimmer2017estimating}, psychologists \citep{bolger2019causal,
  winkelbeiner2019evaluation}, sociologists \citep{xie2012estimating,
  breen2015heterogeneous}, economists
\citep{florens2008identification, djebbari2008heterogeneous}, and
education researchers \citep{morgan2001counterfactuals,
  brand2010benefits}. In short, there is a wide range of fields that
would benefit from a better understanding of ITE. 

% The burst of the
% demand to understand ITE have created a dizzying array of exciting
% applications.

% Traditional approaches for ITE typically start with pre-defined subgroups based on a handful of variables, e.g. gender, age and polygenic risk scores, and then estimate the average effect for each group. With the rise of machine learning over the past decade, there has been growing interests in statistics community to develop new methodologies for ITE. In contrast to traditional approaches, machine learning-based approaches define the subgroups in a data-driven manner or even be fully individualized by treating every subject as a separate group. The data-adaptivity of these methods overcomes at least one clear pitfall of traditional methods that the subgroups defined a priori may not be meaningful for the outcome of interest. With more meaningful subgroups, a finer description of treatment effect heterogeneity is available to improve the evaluation of interventions. 

To move past the average treatment effect as the object of inference,
most existing works have targeted, instead, the {\em conditional
  average treatment effects} (CATE)---although we will introduce a
formal definition later in Section \ref{sec:setup}, this is the
expectation of the ITE conditional on the values of the covariates. % On
% the one hand, 
While CATE naturally provide a richer summary than ATE,
they surely still neglect the inherent variability in the response
(the conditional variance, if you will) which might be crucial for
decision-making; unless, of course, the covariates explain away most
of the variation in ITE. % In practice, however, it is rarely known
% whether the covariates have sufficient explanatory power.
% On the  other hand, if sufficiently many covariates explain away a significant
% fraction of the variability in the response, then uncertainty
% quantification becomes challenging.
As a consequence, two types of variability are of immediate concern:
(1) the variability of the response around the regression function
and, (2) the variability of CATE estimators due to finite
samples. Getting both (1) and (2) under control requires having
sufficiently many covariates to explain away a significant fraction of
the variability in the response, and at the same time, a nearly
perfect estimate of CATE for every value of the covariates. Neither of
these seem to be realistic for everyday causal inference problems.

  In areas like medical science or public policy, point estimates are
  insufficient to inform decisions due to the huge loss potentially
  incurred by wrong actions. A confidence interval, or at least a
  p-value, is required by the U.S.~Food and Drug Administration to
  approve a drug, in order to guarantee sufficient evidence and
  confidence in favor of the drug. The issues are that despite the
  importance of uncertainty quantification, the reliability of modern
  machine learning methods is typically under-studied, and that
  theoretical guarantees are hard to come by. For instance, existing
  theory usually requires strong non-verifiable assumptions and
  asymptotic regimes one does not encounter in practice. This of
  course limits the applications of machine learning methods in
  sensitive causal inference problems.  In addition, we will
  demonstrate later that the confidence intervals for CATE or prediction intervals for ITE
  produced by frequently discussed methods (including Bayesian
  methods) typically have unsatisfactory or unacceptable coverage,
  even in very simple and smooth models with at most ten
  covariates.
% On the other
%   hand, for frequently discussed methods (including Bayesian methods)
%   with uncertainty quantification of either CATE or ITE,  we will
% demonstrate later that their confidence intervals typically have
% unsatisfactory or unacceptable coverage \changed{for both CATE and
%   ITE}, even in very simple models with about a dozen of
% covariates. The lack of uncertainty quantification limits the
% applications of machine learning methods in sensitive causal inference
% problems. 

% Ideally, it is desirable to have interval estimates for ITE
% with guaranteed coverage. \changed{This is more reliable than
%   intervals of CATE, which ignore the inherent variability. }

In this work, we leverage ideas from conformal inference to construct
valid prediction intervals for individual treatment effects under the
potential outcome framework \citep{neyman1923application,
  rubin1974estimating}. In particular, we address two challenges:
\begin{enumerate}[(1)]
\item Construct intervals for ITE with reliable coverage for subjects
  in the study, for which one of the potential oucomes is missing;
\item Construct intervals for ITE with reliable coverage for subjects
  not in the study, and for which both potential oucomes are missing.
\end{enumerate}
To be sure, in the case where a method is able to estimate ITE
accurately, there is no need to distinguish these two tasks. In
practice, however, this is rarely possible due to insufficient sample
sizes, model misspecification, and inherent variability. We will thus see that the second challenge is generally far more ambitious since two potential outcomes are never observed simultaneously; consequently we cannot model them jointly.% We will
% thus see that the second challenge is generally far more ambitious
% since two potential outcomes are never observed simultaneously so that
% one cannot model them jointly.

In brief, Section \ref{sec:counterfactual} shows how the first
challenge reduces to counterfactual inference since one potential
outcome is observed for each subject; here, an interval for the
missing outcome can be shifted into an interval for the individual
treatment effect by contrasting it with the observed outcome.  This
section also introduces methods with guaranteed coverage even under
model misspecification. In particular, for completely randomized or
stratified randomized trials with perfect compliance so that the
propensity score is known but may not be constant, our method
achieves coverage in finite samples \emph{without any assumption other
  than that of operating on ~i.i.d.~samples}. 
% \lihua{By randomized trials I meant either completely randomized
%   trials where the propensity score is a constant or the general
%   blocking experiments where the propensity score is known but varying
%   with $x$.  In the latter case, weighted conformal inference is used
%   and so ~i.i.d.~ samples are required (of course this can be
%   generalized to weighted exchangeable samples but I felt this may be
%   distractive)}. 
For general observational studies under the strong
ignorability assumption \citep{rubin1974estimating}, or randomized
experiments with ignorable compliance, our methods have guaranteed
coverage provided that either the outcome model or the treatment model
is accurately estimated. This is analogous to the \emph{doubly robust}
property applicable to average treatment effect
\citep[e.g.][]{robins1994estimation, kang2007demystifying} which
speaks to the consistency of point
estimates. % , although we emphasize that our double robustness is a novel
% concept for the coverage of interval estimates
 % \ejc{Do we need to say
 %  the last piece?}\lihua{I think it is not needed so I deleted it.}

Having addressed the first challenge, we will see in Section
\ref{sec:effect} that our methods can serve as a stepping stone for
the second. A naive approach, here, would be to apply the
counterfactual inference on both potential outcomes and contrast the
two intervals to induce an interval for ITE.
% Such an approach would
% entirely decouple the pair of potential outcomes and would likely be
% overly conservative. Instead, we introduce a less conservative 
% approach which can exploit any solution to the first task.
Moreover, we introduce another approach 
% \changed{It
%   starts by splitting the data into two folds and applying the counterfactual
%   inference on the second fold,  with the first fold being the training data,  to construct intervals for ITE, which have guaranteed coverage with appropriately chosen
%   counterfactual intervals. It then generalizes and calibrates these intervals
%   to the testing points with both potential outcomes missing. 
% }
%   It
which applies the counterfactual inference to generate intervals for ITE
applicable to subjects in the dataset as an intermediate step, and
trains a model to generalize these intervals to subjects not in the
study.

Finally, our methods easily extend to a widely studied problem which
goes by the name of generalizability, or transportability, or external
validity, and which concerns settings in which there is a
distributional shift between the target population and the study
population \citep[e.g.][]{stuart2011use, tipton2014generalizable}, see
Sections \ref{sec:setup}--\ref{sec:effect}. Furthermore, as explained
in Section \ref{sec:extension}, our methods also naturally adapt to
other causal inference frameworks, such as causal diagrams
\citep{pearl1995causal} and invariant prediction
\citep{peters2016causal}.


%%% Local Variables:
%%% mode: latex
%%% TeX-master: t
%%% End:
